\documentclass[../../main.tex]{subfiles}
\begin{document}

{\bf[解]}

\subparagraph{(1)}

\begin{figure}[htb]
  \centering
  \begin{tikzpicture}[scale=2.6, >=stealth, every node/.style={font=\small}]

    % --- 1. 関数定義・座標パラメータ ---
    % f(x) = x^3/6 + 1/(2x)
    % A: 初期点 (x = 1, y = 2/3 ≈ 0.667)
    % P: (2) で求める点 (x = 2, y = 8/6 + 1/4 = 19/12 ≈ 1.583)
    \coordinate (O) at (0, 0);
    \coordinate (A) at (1.0, {2.0/3.0});
    \coordinate (P) at (2.0, {19.0/12.0});

    % --- 2. 座標軸 ---
    \draw[->, thick] (-0.3, 0) -- (2.8, 0) node[right] {$x$};
    \draw[->, thick] (0, -0.3) -- (0, 2.5) node[above] {$y$};
    \node[below left=2pt] at (O) {$O$};

    % --- 3. 曲線 y = f(x) の描画 ---
    % 全体の曲線（細い実線）
    \draw[thick, gray!50, domain=0.65:2.5, samples=70]
    plot (\x, {\x*\x*\x/6.0 + 1.0/(2.0*\x)});
    \node[gray!70!black, font=\footnotesize, right] at (2.4, {2.4^3/6.0 + 1.0/(4.8)})
    {$y = \frac{x^3}{6} + \frac{1}{2x}$};

    % 点 P が走破した区間 [1, 2] の曲線（太線で強調）
    \draw[line width=2.2pt, blue!85!black, domain=1.0:2.0, samples=50]
    plot (\x, {\x*\x*\x/6.0 + 1.0/(2.0*\x)});

    % --- 4. 走破した弧長 L の表示 ---
    \node[blue!90!black, font=\normalsize, above left=1pt] at (1.45, 1.05) {$L = \frac{t^2}{4}$};

    % --- 5. 始点 A(1, 2/3) と座標投影破線 ---
    \draw[dashed, gray!60] (1.0, 0) -- (A) -- (0, {2.0/3.0});
    \node[below=2pt, font=\scriptsize] at (1.0, 0) {$1$};
    \node[left=2pt, font=\scriptsize]  at (0, {2.0/3.0}) {$\frac{2}{3}$};
    \fill[black] (A) circle (0.8pt) node[below right=1pt] {$A\left(1, \frac{2}{3}\right)$ ($t=0$)};

    % --- 6. 到達点 P(2, 19/12) と座標投影破線 ---
    \draw[dashed, gray!60] (2.0, 0) -- (P) -- (0, {19.0/12.0});
    \node[below=2pt, font=\scriptsize] at (2.0, 0) {$2$};
    \node[left=2pt, font=\scriptsize]  at (0, {19.0/12.0}) {$\frac{19}{12}$};
    \fill[red!85!black] (P) circle (0.9pt) node[below right=2pt, red!90!black] {$P(x(t), y(t))$};

    % --- 7. 点 P における速度ベクトル v (接線方向) ---
    % 傾き dy/dx = 15/8 = 1.875
    % 方向角: atan(15/8) ≈ 61.93°
    % ベクトルの長さ = 0.75
    \pgfmathsetmacro{\angTan}{atan(15.0/8.0)}
    \coordinate (V_end) at ($ (P) + ({\angTan}:0.75) $);

    % 接線（延長破線）
    \draw[dashed, teal!70!black] ($ (P) - ({\angTan}:0.5) $) -- (P);
    % 速度ベクトル矢印
    \draw[->, line width=1.8pt, teal!85!black] (P) -- (V_end);
    %  node[above right=1pt, font=\normalsize] {$\vec{v} \parallel \begin{pmatrix}8\\15\end{pmatrix}$};

    % % 傾き比率を示す直角三角形 (8 : 15)
    % \coordinate (T_dx) at ($ (P) + (0.35, 0) $);
    % \coordinate (T_dy) at ($ (P) + (0.35, {0.35*15.0/8.0}) $);
    % \draw[thin, teal!70!black] (P) -- (T_dx) -- (T_dy);
    % \node[below, font=\tiny, teal!90!black] at ($ (P)!0.5!(T_dx) $) {$8$};
    % \node[right, font=\tiny, teal!90!black] at ($ (T_dx)!0.5!(T_dy) $) {$15$};

  \end{tikzpicture}
  \caption{曲線 $y=\dfrac{x^3}{6}+\dfrac{1}{2x}$ 上を運動する点 $P$ と速度ベクトル $\vec{v}$}
  \label{fig:curve_motion_tangent}
\end{figure}

題意の曲線
\begin{align}
  y = f(x) = \dfrac{x^3}{6} + \dfrac{1}{2x} \label{eq:1}
\end{align}
について考える．

時刻$s$での点$P$の速度の大きさは$\dfrac{s}{2}$だから，点$P$は$[0,t]$の間に$y=f(x)$上を
\begin{align}
  \int_0^t\dfrac{s}{2}ds=\dfrac{t^2}{4} \label{eq:2}
\end{align}
だけ進む．これが，曲線上で$(1,2/3)$から$(x,y)$までの長さ$L$に等しい．
ここで
\begin{align}
  \dfrac{dy}{dx} = f'(x) = 2x^2 - \dfrac{1}{2x^2} \label{eq:3}
\end{align}
だから具体的に$L$を計算すると
\begin{align}
  L
  &=\int_{1}^{x}\sqrt{1+(dy/dx)^2}dx \nonumber\\
  &=\int_{1}^{x}\sqrt{1+\left(2x^2 - \dfrac{1}{2x^2}\right)^2}dx \nonumber\\
  &=\int_{1}^{x}\sqrt{1+4x^4 + \dfrac{1}{4x^4} - \dfrac{1}{2}}dx \nonumber\\
  &=\int_{1}^{x}\sqrt{\dfrac{(x^2+1/x^2)^2}{4}}dx \nonumber\\
  &=\dfrac{1}{2}\int_{1}^{x}\left(x^2+\dfrac{1}{x^2}\right) dx \nonumber\\
  &=\dfrac{1}{2}\left[\dfrac{x^3}{3}-\dfrac{1}{x}\right]_{1}^{x} \nonumber\\
  &=\dfrac{1}{2}\left(\dfrac{x^3}{3}-\dfrac{1}{x}+\dfrac{2}{3}\right) \label{eq:4}
\end{align}
である．\cref{eq:2,eq:4}が等しいので，求める$t$と$x$の関係は
\begin{align}
  t^2=2\left(\dfrac{x^3}{3}-\dfrac{1}{x}+\dfrac{2}{3}\right) \label{eq:5}
\end{align}
である．$\cdots$(答)

\subparagraph{(2)}

$\vec{v}$は$P$での曲線の接線方向のベクトルと平行である．
\cref{eq:3}から，$(1,f(x))\parallel(8,15)$の時，
\begin{align}
  &f(x)=\dfrac{15}{8} \\
  &\left(x^2-\dfrac{1}{x^2}\right)=\dfrac{15}{4} \quad (\because \text{\cref{eq:3}})
\end{align}
である．両辺$x^2$倍して，$p=x^2$とおくと
\begin{align}
  &x^4 - \dfrac{15}{4}x^2 - 1 = 0 \\
  &p^2 - \dfrac{15}{4}p - 1 = 0 \\
  &(4p+1)(p-4)=0 \\
  &p=4,\dfrac{-1}{4}
\end{align}
である．題意から$x>1$より$p=x^2>1$だからこれを満たすのは$p=4$つまり$x=2$の時である．
この時の$t$は\cref{eq:5}に$x=2$を代入して
\begin{align}
  t^2=2\left(\dfrac{2^3}{3}-\dfrac{1}{2}+\dfrac{2}{3}\right)=\dfrac{17}{3}
  \quad \therefore
  t = \sqrt{17/3}
\end{align}
である．$\cdots$(答)

\newpage
\end{document}
